%%%%%%%%%%%%%%%%%%%%%%%%%%%%%%%%%%%%%%%%%%%%%%%%%%%%%%%%%%%%%%%%%%%%%%%%%%%%%%%%%
%	CONCLUSION AND FUTURE RESEARCH DIRECTIONS
%   Mandated by Project Task 4.2 ("future research directions").
%%%%%%%%%%%%%%%%%%%%%%%%%%%%%%%%%%%%%%%%%%%%%%%%%%%%%%%%%%%%%%%%%%%%%%%%%%%%%%%%%
\lhead{Conclusion and Future Work}
\chapter{Conclusion and Future Research Directions}
\label{SecConclusion}

\section{Summary of Findings}
\label{Conc:Summary}

We investigated a Genetic Algorithm for the Traveling Salesman Problem on the
kroA100 benchmark (100 cities, known optimal tour length 21\,282). The primary
pipeline pairs naive single-point crossover with a random-insertion repair
step. We compared it against two alternatives: a penalty function and the
feasibility-preserving PMX operator.

From the three-arm comparison (\sref{T4:Boxplots}), we found that repair and
PMX performed similarly on kroA100, with mean best fitness values of $49\,778$
and $50\,655$ respectively, while the penalty approach was clearly worse at
$57\,948$. This suggests that naive crossover with repair is a viable
alternative to feasibility-preserving operators like PMX.

The $1\,620$-run parameter sweep (\sref{T3:Sensitivity}) showed that selection
method matters the most, with tournament selection outperforming roulette
across all settings. Mutation rate was the next most important factor, as
higher rates helped avoid premature convergence. Population size and crossover
rate had smaller effects. The best configuration we found used population size
$200$, crossover rate $0.95$, mutation rate $0.10$ with tournament selection,
achieving a mean of $49\,347$.

\section{Limitations}
\label{Conc:Limitations}

\begin{itemize}[nosep]
    \item All experiments were run on kroA100 only, so the results may not
    generalise to larger instances or different TSP variants.

    \item The $\sim\!100\%$ gap from the known optimum is expected since we
    did not use any local search like 2-opt or Lin--Kernighan to refine
    the tours.
\end{itemize}

\section{Future Research Directions}
\label{Conc:FutureWork}

\begin{enumerate}[nosep]
    \item Add 2-opt or Lin--Kernighan local search after crossover and
    repair to reduce the optimality gap (\cite{liu2014powerfulga}).

    \item Test on larger benchmarks like kroA200 and kroA500 to see if
    the repair versus PMX versus penalty ranking still holds.

    \item Try adaptive penalty weights using the violation-factor method
    from \cite{chehouri2016violationfactor} to check if the penalty
    approach could work better with proper tuning.
\end{enumerate}
